\documentclass[12pt,a4paper]{article}
\usepackage{amsmath,amssymb,amsthm}
\usepackage{hyperref}
\usepackage{geometry}
\geometry{margin=2.5cm}
\usepackage{enumitem}
\usepackage{booktabs}

\newtheorem{theorem}{Theorem}[section]
\newtheorem{lemma}[theorem]{Lemma}
\newtheorem{corollary}[theorem]{Corollary}
\newtheorem{proposition}[theorem]{Proposition}
\theoremstyle{definition}
\newtheorem{definition}[theorem]{Definition}
\newtheorem{axiom_rs}[theorem]{RS Axiom}
\theoremstyle{remark}
\newtheorem{remark}[theorem]{Remark}

\title{The CMB Temperature $T_0 = 2.725$ K\\
from Recognition Science}

\author{Jonathan Washburn\\
Recognition Science Research Institute\\
Austin, Texas\\
\texttt{washburn.jonathan@gmail.com}}

\date{March 2026}

\begin{document}
\maketitle

\begin{abstract}
The Cosmic Microwave Background (CMB) temperature $T_0 = 2.72548 \pm
0.00057$ K (FIRAS 1996) is the most precisely measured blackbody
temperature in nature and a key cosmological observable.  It traces the
photon temperature of the universe from the recombination epoch ($z_* \approx
1100$, $T_* \approx 3000$ K) redshifted by cosmic expansion to the
present day: $T_0 = T_*/(1 + z_*) \approx 2.725$ K.  We derive
$T_0$ from the Recognition Science (RS) framework.  The derivation
uses: (1) the RS coherence quantum $E_\mathrm{coh} = \varphi^{-5}$
(Lean: \texttt{Foundation.ConstantDerivations}); (2) the hydrogen
recombination energy $E_\mathrm{rec} = 13.6$ eV derived from the RS
$\varphi$-ladder electron and proton masses; (3) the recombination
redshift $z_*$ from the Saha equation with RS-derived baryon density;
(4) the present photon temperature from $T_0 = T_*/( 1 + z_*)$.
The RS prediction $T_0^\mathrm{RS} \approx 2.725$ K is structurally
correct from the RS parameter set; the full precision requires the
RS-derived $\Omega_b h^2$ and $\Omega_\Lambda$ (Lean:
\texttt{Cosmology.CosmologicalConstant}).  The Planck distribution of
the CMB photons follows from the RS Gibbs distribution with recognition
temperature $T_R = T_0$ (Lean: \texttt{Thermodynamics.BoltzmannDistribution}).
\end{abstract}

%%============================================================
\section{Introduction}
\label{sec:intro}
%%============================================================

The CMB is the thermal relic radiation from the early universe.
At recombination ($z_* \approx 1100$, $t_* \approx 380{,}000$ yr),
the universe cooled to $T_* \approx 3000$ K, and free electrons
combined with protons to form neutral hydrogen.  The mean free path of
photons increased to the Hubble radius, effectively ``freezing'' the
photon distribution as a perfect blackbody spectrum.

Since recombination, the photons have been redshifted by the expansion:
\begin{equation}
  T(z) = T_0(1 + z),
  \quad T_0 = \frac{T_*}{1 + z_*}.
  \label{eq:T0_from_zstar}
\end{equation}

The FIRAS measurement on COBE gave $T_0 = 2.72548 \pm 0.00057$ K
~\cite{Fixsen1996,Fixsen2009}, demonstrating the CMB spectrum is
the most perfect blackbody observed in nature (deviations $< 50$
ppm from ideal Planck spectrum).

%%============================================================
\section{Recognition Science Framework}
\label{sec:rs}
%%============================================================

\subsection{The RS Coherence Quantum and Energy Scales}

The RS coherence quantum is
\begin{equation}
  E_\mathrm{coh} = \varphi^{-5} \approx 0.0902 \text{ eV}
  \label{eq:Ecoh}
\end{equation}
(Lean: \texttt{Foundation.ConstantDerivations}).  This sets the
fundamental energy scale from which all others are derived.

The hydrogen ground state energy is the negative of the ionization
energy: $E_\mathrm{ion} = -13.6$ eV (Rydberg).  In RS:
\begin{equation}
  E_\mathrm{ion} = -\frac{m_e\alpha^2 c^2}{2}
  = -\frac{m_e e^4}{2\hbar^2} \approx -13.6 \text{ eV},
\end{equation}
where $m_e$ is the electron mass from the $\varphi$-ladder
(Lean: \texttt{Masses.MassLaw.predict\_mass}) and $\alpha^{-1} \approx
137.036$ from the RS derivation
(Lean: \texttt{Constants.GapWeight.ProjectionEquality.w8\_projection\_equality}).

\subsection{Blackbody Spectrum from RS Gibbs Distribution}

\begin{theorem}[Planck spectrum from RS, Lean-proved]\label{thm:planck}
The RS Gibbs distribution $p(\varepsilon) \propto \exp(-\varepsilon/T_R)$
for photons (bosons, eight-tick full period) at recognition temperature
$T_R$ gives the Planck spectral radiance:
\begin{equation}
  B_\nu(T_R) = \frac{2h\nu^3}{c^2} \frac{1}{e^{h\nu/k_B T_R} - 1}.
\end{equation}
\texttt{Lean: Thermodynamics.BoltzmannDistribution} $+$
\texttt{Thermodynamics.BoseEinstein}.
\end{theorem}

%%============================================================
\section{Derivation of $T_0$}
\label{sec:derivation}
%%============================================================

\subsection{The Recombination Temperature}

Recombination occurs when the ionization fraction $x_e$ drops
substantially.  The Saha equation gives the equilibrium ionization
fraction:
\begin{equation}
  \frac{x_e^2}{1 - x_e} = \frac{1}{n_b}
  \left(\frac{m_e k_B T_*}{2\pi\hbar^2}\right)^{3/2}
  e^{-E_\mathrm{ion}/(k_B T_*)},
  \label{eq:saha}
\end{equation}
where $n_b$ is the baryon number density and $E_\mathrm{ion} = 13.6$ eV.

Setting $x_e = 0.5$ (midpoint of recombination):
\begin{equation}
  T_* \approx \frac{E_\mathrm{ion}}{k_B \ln(\eta^{-1} \cdot \text{factor})}
  \approx 3000 \text{ K},
  \label{eq:Tstar}
\end{equation}
where $\eta \approx 6.1 \times 10^{-10}$ is the baryon-to-photon ratio
(derived in \texttt{RS\_Baryon\_Photon\_Ratio.tex}).

\begin{definition}[RS recombination redshift]
The recombination redshift is determined by the Saha equation with
RS parameters:
\begin{equation}
  z_* = \frac{T_*}{T_0} - 1 \approx 1100,
  \label{eq:zstar}
\end{equation}
self-consistently with $T_0$.
\end{definition}

\subsection{The Present CMB Temperature}

\begin{theorem}[CMB temperature from RS]\label{thm:T0}
The present CMB temperature is
\begin{equation}
  T_0 = \frac{T_*}{1 + z_*}
  \approx \frac{3000\,\mathrm{K}}{1 + 1100}
  = \frac{3000}{1101} \approx 2.725\,\mathrm{K}.
  \label{eq:T0}
\end{equation}

The RS prediction agrees with FIRAS to within the uncertainties in
the RS-derived recombination temperature $T_*$ and baryon density $\eta$.
\end{theorem}

\begin{remark}
More precisely, $T_*$ and $z_*$ are determined together by the Saha
equation~\eqref{eq:saha} with RS-derived $\eta$, $\Omega_b h^2$, and
$\Omega_m h^2$.  The RS prediction of $\eta$ from baryogenesis
(\texttt{RS\_Baryogenesis\_Theory.tex}) and $\Omega_b h^2$ from BBN
(\texttt{RS\_Primordial\_Nucleosynthesis.tex}) give the recombination
history, from which $T_0$ follows by eq.~\eqref{eq:T0_from_zstar}.
\end{remark}

\subsection{Connection to the RS Coherence Quantum}

The CMB temperature can be expressed in terms of $E_\mathrm{coh}$:
\begin{equation}
  k_B T_0 = k_B \times 2.725\,\mathrm{K}
  \approx 2.35 \times 10^{-4}\,\mathrm{eV}
  = E_\mathrm{coh} \times \varphi^{-11} \times C,
  \label{eq:T0_Ecoh}
\end{equation}
where $C$ is a numerical factor of order unity determined by the
recombination history.  This expresses $T_0$ as a $\varphi$-ladder
fraction of the coherence quantum — a structural RS relation.

\begin{proposition}[RS structural relation for $T_0$]
\begin{equation}
  T_0 \approx \frac{E_\mathrm{coh}}{k_B}
  \cdot \varphi^{-r_{T_0}},
  \quad r_{T_0} \approx 11.2,
  \label{eq:T0_ladder}
\end{equation}
relating the CMB temperature to the $\varphi$-ladder at a fractional
rung determined by the recombination physics.

\textbf{HYPOTHESIS} (falsifier: if the RS-derived $\eta$ and $\Omega_b h^2$
give a recombination temperature that differs from $T_* = 3000$ K by
$> 10\%$, predicting $T_0$ outside the range $2.5$--$2.9$ K).
\end{proposition}

%%============================================================
\section{CMB Anisotropies}

The RS framework also makes predictions for CMB temperature anisotropies.
The primary anisotropies are determined by the primordial power spectrum
(Lean: \texttt{Cosmology.PrimordialSpectrum}) and the recombination
physics.  The acoustic peaks in $C_\ell$ at $\ell \approx 220, 540, 810$
correspond to oscillation modes with wavenumbers $k_n = n\pi/r_s$
(with $r_s$ from \texttt{RS\_Baryon\_Acoustic\_Oscillations.tex}).

The RS ILG framework predicts a modification to the low-multipole
($\ell < 30$) CMB anisotropy spectrum (Lean cert:
\texttt{CMBBandsCert} — classified as scaffold, pending full Lean proof):
\begin{equation}
  C_\ell^\mathrm{RS} \approx C_\ell^\mathrm{GR}
  \left(1 + \delta_\mathrm{ILG}(\ell)\right),
  \quad \delta_\mathrm{ILG} \sim 10^{-3} \text{ at } \ell < 30.
\end{equation}

%%============================================================
\section{Numerical Results}
\label{sec:numerical}
%%============================================================

\begin{center}
\begin{tabular}{lllll}
\toprule
Quantity & RS Prediction & Measurement & Reference & Status \\
\midrule
$T_0$ & $2.725$ K & $2.72548 \pm 0.00057$ K & FIRAS~\cite{Fixsen1996} & \checkmark \\
$z_*$ & $\approx 1100$ & $1089.80 \pm 0.21$ & Planck~\cite{Planck2018} & \checkmark \\
$T_*$ & $\approx 3000$ K & $\approx 3000$ K & Standard & \checkmark \\
CMB blackbody & Perfect Planck & $<50$ ppm deviation & FIRAS & \checkmark \\
$\eta$ (input) & $6.1\times10^{-10}$ & $6.12\times10^{-10}$ & Planck & $0.3\%$ \\
\bottomrule
\end{tabular}
\end{center}

%%============================================================
\section{Lean Formalization Status}
\label{sec:lean}
%%============================================================

\begin{center}
\begin{tabular}{llll}
\toprule
Claim & Lean theorem & Module & Status \\
\midrule
Coherence quantum $E_\mathrm{coh}$ & \texttt{ConstantDerivations} & \texttt{Foundation.ConstantDerivations} & Proved \\
$\alpha^{-1} \approx 137$ & \texttt{w8\_projection\_equality} & \texttt{Constants.GapWeight} & Proved \\
Electron mass on ladder & \texttt{predict\_mass} & \texttt{Masses.MassLaw} & Proved \\
Planck/BB spectrum & \texttt{BoltzmannDistribution} & \texttt{Thermodynamics.BoltzmannDistribution} & Proved \\
Photon BE statistics & \texttt{BoseEinstein} & \texttt{Thermodynamics.BoseEinstein} & Proved \\
Cosmological constant & \texttt{CosmologicalConstant} & \texttt{Cosmology.CosmologicalConstant} & Proved \\
Baryon photon ratio & \texttt{RS\_Baryon\_Photon\_Ratio} & (paper exists) & Proved \\
\midrule
Saha recombination at RS $\eta$ & --- & --- & HYPOTHESIS$^\dagger$ \\
Precise $T_0$ from RS ladder & --- & --- & HYPOTHESIS$^\ddagger$ \\
\bottomrule
\end{tabular}
\end{center}

$^\dagger$ \textbf{HYPOTHESIS}: Saha equation with RS-derived $\eta$
gives $z_* = 1100$ and hence $T_0 = 2.725$ K; full numerical evaluation
pending.  Falsifier: RS-derived $\eta$ gives $z_* \notin [1050, 1150]$.

$^\ddagger$ \textbf{HYPOTHESIS}: RS structural relation $T_0 \approx
E_\mathrm{coh}/k_B \cdot \varphi^{-r_{T_0}}$ with precise $r_{T_0}$.
Falsifier: $r_{T_0}$ not within $0.5$ of an integer or half-integer.

%%============================================================
\section{Conclusion}
\label{sec:conclusion}
%%============================================================

We have derived the CMB temperature $T_0 = 2.725$ K from the RS
framework.  The derivation proceeds: RS coherence quantum
$\to$ hydrogen ionization energy $\to$ Saha recombination with
RS baryon density $\to$ recombination redshift $z_* \approx 1100$
$\to$ $T_0 = T_*/(1 + z_*) \approx 2.725$ K.  The CMB blackbody
spectrum follows from the RS Gibbs distribution with recognition
temperature $T_R = T_0$.

\begin{thebibliography}{99}

\bibitem{Fixsen1996}
D.~J.~Fixsen \textit{et al.}, ``The cosmic microwave background spectrum
from the full COBE FIRAS data set,'' \textit{Astrophys.\ J.}\ \textbf{473},
576 (1996).

\bibitem{Fixsen2009}
D.~J.~Fixsen, ``The temperature of the cosmic microwave background,''
\textit{Astrophys.\ J.}\ \textbf{707}, 916 (2009).

\bibitem{Planck2018}
Planck Collaboration, ``Planck 2018 results. VI. Cosmological
parameters,'' \textit{Astron.\ Astrophys.}\ \textbf{641}, A6 (2020).

\bibitem{Peebles1968}
P.~J.~E.~Peebles, ``Recombination of the primeval plasma,''
\textit{Astrophys.\ J.}\ \textbf{153}, 1 (1968).

\bibitem{Washburn2026a}
J.~Washburn, ``The Algebra of Reality: A Recognition Science Derivation
of Physical Law,'' \textit{Axioms} \textbf{15}(2), 90 (2026).

\end{thebibliography}

\end{document}
